\documentstyle[% 


righttag% 


,11pt% 


,amssymb% 


,russian%               remove in Eng. 


,izvru%                 Set izven% in Eng. 


% ,izvru-ks             for brief communications 


]{amsart} 


 


% 


 


\newcommand{\updef}{\overset{\mathrm{def}}{=}} 


\newcommand{\const}{\operatorname{const}} 


\newcommand{\var}{\operatorname{var}} 


\newcommand{\tts}[1]{} 


 


\theoremstyle{definition} 


\theorembodyfont{\rm} 


\newtheorem{notenonum}{\hskip\parindent Замечание} 


\renewcommand{\thenotenonum}{} 


 


%\hoffset=-15mm%                  For Eng. 


\setcounter{page}{49} 


\god{2000} 


\nomer{10~(461)} %                        Remove in Eng. 


%\nomer{Vol.\,44, No.\,10}%               Set in Eng. 


%\str{\pageref{begin}--\pageref{end}}%  Set in Eng. 


\udk{517.538} 


 


\author{\it Ф.Г.\,Насибов} 


% NB:   ^^ remove in Eng. 


\title{Об Экстремальных задачах в классах целых  


функций, ограниченных на вещественной оси} 


 


\date{} 


%\pagestyle{myheadings}%         Set in Eng. 


\pagestyle{plain} %              Remove in Eng. 


\begin{document} 


%\thispagestyle{plain}%          Set in Eng. 


\maketitle 


\label{begin} 


 


\section{Введение} 


 


{\bf 1.} Следуя ([1], с.\,182), обозначим через $B_\sigma$ класс целых 


функций 


конечной степени $\le\sigma$, ограниченных на оси $R=(-\infty,\infty)$. 


С.Н.\,Бернштейн впервые ввел этот класс [2] и доказал, что $\forall f\in 


B_\sigma$ имеет место неулучшаемое неравенство 


$$ 


|f'(x)|\le\sigma\|f\|_C=\sigma\sup_{x\in R}|f(x)| 


\quad \forall x\in R. 


$$ 


 


В дальнейшем появились многочисленные исследования, в которых это 


неравенство обобщалось в различных направлениях, переносилось на другие 


классы функций с разнообразными применениями ([1], [3]--[5]). 


 


В данной работе сделана попытка рассмотреть с единой точки зрения 


все такие частные задачи и показать, что все они укладываются в единую 


схему (с частичным анонсом в [6]). В классах целых функций 


$W_{\sigma,2}$ (порядка 


$\rho=1$ и принадлежащих к $L_2$), $D_{\sigma,1/2}$ и $B_{\sigma,1/2}$ 


($D$, $B$ --- классы целых функций 


порядка $\rho=1/2$) такие общие исследования проводились автором ранее 


([7]--[10]). Исследования в равномерной метрике осложняются отсутствием 


общего вида линейного функционала в пространстве $C(-\infty,\infty)$. 


 


К экстремальным задачам в различных классах аналитических функций 


такой общий подход сделан на основе теоремы Хана--Банаха давно, начиная 


с 1949 г. [11], поставлены и исследованы достаточно общие экстремальные 


задачи в различных классах аналитических функций. 


 


Основные результаты здесь принадлежат С.Я.\,Хавинсону, Г.Ц.\,Тумаркину, 


В.\,Рогозинскому, А.\,Макинтайру, Г.\,Шапиро и др. Обзор результатов в 


этом направлении можно найти в статьях [12]--[14], которым 


предшествовали 


основополагающие работы М.Г.\,Крейна [15] и С.М.\,Ни\-коль\-с\-ко\-го [16]. 


Отметим 


также работу Е.А.\,Горина [17], где неравенство С.Н.\,Бернштейна 


исследуется с точки зрения теории операторов. 


 


Наш подход основан на описании аннуляторов некоторых классов целых 


функций, которые позволяют устанавливать соотношения двойственности, 


выявлять ряд особенностей исследуемых задач, устанавливать связь 


экстремальных задач с проблемой наилучшего приближения. 


\medskip 


 


{\bf 2.} Введем обозначения и определения классов, которые являются 


объектом наших исследований [6]. 


 


1) Через $H_\sigma$ обозначаем класс целых функций степени $\le\sigma$. 


 


2) $V[-\sigma,\sigma]\overset{\mathrm{def}}{\equiv}


\Big\{\nu:\var(\nu:[-\sigma,\sigma])\equiv 


\overset{\sigma}{\underset{-\sigma}{V}}(\nu)= 


\int\limits^\sigma_{-\sigma}|d\nu|<+\infty\Big\}$ --- класс мер 


$\nu(t)$ 


с ограниченным изменением на $[-\sigma,\sigma]$. 


 


В случае $V(-\infty,\infty)$ будем употреблять обозначениe $V\equiv V_R$. 


 


3) $\cal M_\sigma\updef\Big\{f\in 


H_\sigma:f(z)=\int\limits^\sigma_{-\sigma} e^{izt}\varphi(t)dt 


\ \forall\varphi\in L_1(-\sigma,\sigma)\Big\}$. 


 


4) $\cal N_\sigma\updef\Big\{f\in 


H_\sigma:f(z)=\int\limits^\sigma_{-\sigma} e^{izt}d\nu(t)\ \forall 


\nu\in V(-\sigma,\sigma)\Big\}$. 


 


5) $\cal N^+_\sigma\updef\{f\in\cal N_\sigma:\nu\uparrow\}$. 


 


6) $\cal N_{\sigma,0}\updef\{f\in\cal N_\sigma:f(x)\to0$ при 


$|x|\to\infty\}$. 


 


Отметим очевидные вложения 


$$ 


W_{\sigma,2}\subset\cal M_\sigma\subset 


\cal N_{\sigma,0}\subset\cal N_\sigma\subset B_\sigma. 


$$ 


 


Здесь нетривиальным является то, что $\cal M_\sigma\not\equiv\cal 


N_{\sigma,0}$, а точнее, $\cal N_{\sigma,0}\setminus\cal M_\sigma$ 


не пусто. Для класса $\cal N_{\sigma,0}$ можно утверждать, что меры 


$\nu(t)$, входящие в 


интегральные представления функций $f\in\cal N_{\sigma,0}$, являются 


непрерывными функциями ([18], с.\,35).





Величина 


$$ 


L=\underset{|x|\to\infty}{\varlimsup}|f(x)|\quad 


\bigg(f(x)=\int^\sigma_{-\sigma}e^{ixt}d\nu(t) \bigg) 


$$ 


в принципе может быть любым числом между 0 и 1. 


Известны примеры, когда $L=0$, $L=1$. Сингулярные меры, для которых 


$L$ равно заданному числу из отрезка $[0,1]$, были построены Шварцем еще 


в 1941 г. ([18], с.\,35). Следовательно, для некоторой сингулярной меры 


$\nu$ 


\; $f(x)\to0$ ($|x|\to\infty$). Это показывает, что существует функция 


$f_0(x)\in\cal N_{\sigma,0}$, 


которая не входит в $\cal M_\sigma$. 


 


В свое время на семинаре в МИАН проф. С.Б.\,Стечкин указал на то, что 


этот факт можно вывести еще из ранних работ Д.Е.\,Меньшова о


нуль-тригонометрических рядах. 


 


Теперь установим два утверждения, которые понадобятся далее. 


 


\begin{lem} 


$\forall f(z)\in B_\sigma$ существует последовательность 


$\{g_n(z)\}\subset\cal N_\sigma$, 


удовлетворяющая следующим условиям$:$ 


 


$1)$ $\{g_n(z)\}$ равномерно ограничена на $R=(-\infty,\infty);$ 


 


$2)$ $f(x)=\lim\limits_{n\to\infty}g_n(x)$ $\forall x\in R$, причем 


сходимость равномерная на каждом конечном 


отрезке оси $R$. 


\end{lem} 


 


\begin{pf} 


Пусть $\Lambda_n(f;x)$ --- полиномы Левитана функции $f\in B_\sigma$, 


\begin{gather*} 


\Lambda_n(f;x)=\sum^n_{k=-n}\frac{\sigma}{n}E_{\sigma/n} 


\bigg(\frac{k\sigma}{n}\bigg)\exp\bigg\{i\frac{k\sigma}{n}x \bigg\},\\ 


% 


E_\delta(x)=E_\delta(f;x)=\frac{1}{2\pi}\int^\infty_{-\infty} 


e^{-ixt}\bigg(\frac{\sin\delta t/2}{t/2} \bigg)^2 


f(t)dt, 


\end{gather*} 


а функция $\nu_n(t)$ с ограниченным изменением на $[-\sigma,\sigma]$ 


делает скачки 


$\lambda_k=\frac{\sigma}{n}E_{\sigma/n}(t_k)$ в точках 


$t_k=\frac{k\sigma}{n}$ ($\pm k=0,1,\dotsc,n$) и эти $t_k$ --- 


единственные 


точки роста $\nu_n(t)$. Тогда получаем


$$ 


g_n(x)=\int^\sigma_{-\sigma}e^{ixt}d\nu_n(t)= 


\sum^n_{k=-n}\frac{\sigma}{n}E_{\sigma/n}(t_k) 


\exp(ixt_k)\equiv\Lambda_n(f;x). 


$$ 


Остается учесть теорему Б.М.\,Левитана ([17], с.\,193), что приводит к 


утверждению 2). Справедливость же 1) следует из того свойства полиномов 


$\Lambda_n(f;x)$, что если на $R$\; $|f(x)|\le M$, то 


$|\Lambda_n(f;x)|\le M$. 


\end{pf} 


 


{\bf Замечание.} В книге [19] Р.\,Боаса имеется утверждение 2), но для 


нас важным добавлением является п.\,1) (у Боаса другой метод 


доказательства). 


 


\begin{lem} 


Класс $\cal N^+_\sigma$ замкнут относительно равномерной сходимости на 


$R=(-\infty,\infty)$. 


\end{lem} 


 


{\bf Доказательство.} В самом деле, пусть последовательность 


$\{f_n(x)\}\subset\cal N^+_\sigma$ 


и $\{\nu_n(t)\}$ --- соответствующая ей последовательность мер по 


формуле 


$$ 


f_n(x)=\int^\sigma_{-\sigma}e^{ixt}d\nu_n(t), 


\quad \nu_n(t)\uparrow\ \; (\forall n). 


$$ 


Если $f_n(x)\to f(x)$ равномерно на $R$, то $\{\nu_n(t)\}$ слабо 


сходится к мере $\nu(t)$ 


([18], сс.\,62, 73), причем $\nu(t)\uparrow$ и 


$$ 


f(x)=\int^\sigma_{-\sigma}e^{ixt}d\nu(t)\in\cal N^+_\sigma. 


\quad\square 


$$ 


 


\section{Описание аннуляторов и некоторые их свойства} 


 


{\bf 1.} Пусть $E$ --- некоторое линейное пространство и 


$F\subset E$, $E^*$ --- сопряженное с $E$ пространство. 


 


Множество всех линейных функционалов $l\in E^*$, каждый из которых 


обращается в нуль на каждом элементе $x\in F$, называется 


{\it аннулятором} 


класса $F$ и обозначается через $F^\bot$. 


 


\begin{thm}[{\series{m}\shape{n}\selectfont [6]}] 


Аннуляторы классов $\cal M_\sigma$ и $\cal N_{\sigma,0}$ совпадают и 


каждый из 


них состоит из мер $\mu(t)$ с ограниченным изменением на $R$, 


преобразование 


Фурье--Стильтьеса которых равно нулю на отрезке $[-\sigma,\sigma]$ $($и 


только из 


таких функций$):$ 


$$ 


\cal M^\bot_\sigma\equiv\cal N^\bot_{\sigma,0}=\bigg\{ 


\mu\in V:\wh\mu(x)=\int^\infty_{-\infty} 


e^{ixt}d\mu(t)=0\quad\text{при}\quad |x|\le\sigma\bigg\}. 


$$ 


\end{thm} 


 


\begin{pf} 


Вначале рассмотрим класс $\cal M_\sigma$. Из определения этого 


класса следует, что $\forall f\in\cal M_\sigma:f(x)\to0$ 


($|x|\to\infty$) (теорема Римана--Лебега). 


Следовательно, $\cal M_\sigma$ является подпространством пространства 


$C_0$-функций, 


непрерывных на $R$ и стремящихся к нулю при $|x|\to\infty$. По теореме 


М.Г.\,Крейна ([4], с.\,123) всякий линейный функционал $l$ над $C_0$ 


представим 


в виде 


$$ 


l(f)=\int_R f(t)d\mu(t),\quad \mu\in V_r,\quad 


\|l\|=\int_R|d\mu|. 


$$ 


Тогда получим, что $\cal M^\bot_\sigma$ состоит из тех и только тех мер 


$\mu\in V_R$, для 


которых $l(f)=0$ ($\forall f\in\cal M_\sigma$). 


 


Пользуясь определением класса $\cal M_\sigma$, приходим к заключению, 


что 


$$ 


l(f)=\int^\sigma_{-\sigma}\varphi(x)\wh\mu(x)dx,\quad 


\wh\mu(x)=\int_R e^{ixt}d\mu(t). 


$$ 


Поэтому $l(f)=0$ тогда и только тогда, когда $\wh\mu(x)=0$ при 


$-\sigma\le x\le\sigma$. 


 


Теперь покажем, что $\cal N^\bot_{\sigma,0}\equiv\cal M^\bot_\sigma$. 


Пусть $\mu(t)\in\cal M^\bot_\sigma$ и $f\in\cal N_{\sigma,0}$. В силу 


представления $f\in\cal N_{\sigma,0}$ находим 


$$ 


l(f)=\int_R f(t)d\mu(t)=\int^\sigma_{-\sigma} 


\wh\mu(x)d\nu(x)=0\quad (\forall f\in\cal N_{\sigma,0}), 


$$ 


т.\,е. $\mu\in\cal N^\bot_{\sigma,0}$. Таким образом, $\cal 


M^\bot_\sigma\subset\cal N^\bot_{\sigma,0}$. Обратное вложение $\cal 


N^\bot_{\sigma,0}\subset\cal M^\bot_\sigma$ 


является следствием вложения $\cal M_\sigma\subset\cal N_{\sigma,0}$. 


\end{pf} 


 


{\bf 2.} Теперь рассмотрим самый широкий класс $B_\sigma\subset 


C_M(R)$, причем $C_M(R)$ 


обозначает пространство всех непрерывных и ограниченных на $R$ функций. 


 


Так как общий вид линейного функционала над $C_M(R)$ неизвестен, то 


мы не можем дать полное описание аннулятора $B^\bot_\sigma$. Поэтому, 


несколько 


отступая от общего определения аннулятора $B^\bot_\sigma$ (также $\cal 


N^\bot_\sigma$), будем 


подразумевать под $B^\bot_\sigma$ множество функционалов вида 


$$ 


l(f)=\int_R f(t)d\mu(t),\quad \mu\in V_R, 


$$ 


обращающихся в нуль на каждом элементе из $B_\sigma(N_\sigma)$. 


(Возможно, в $C_M(R)$ 


существуют и другие функционалы, аннулирующие $B_\sigma(N_\sigma)$.) С 


таким уточнением 


имеет место 


 


\begin{thm} 


Справедливы следующие равенства$:$ 


$$ 


B^\bot_\sigma=\cal N^\bot_\sigma=\cal M^\bot_\sigma=\cal 


N^\bot_{\sigma,0}. 


$$ 


\end{thm} 


 


\begin{pf} 


С одной стороны, ясно, что $B^\bot_\sigma\subset\cal 


N^\bot_\sigma\subset\cal N^\bot_{\sigma,0}$. С другой 


стороны, $\forall\mu\in\cal N^\bot_{\sigma,0}$ и $\forall f\in 


B_\sigma$ (в силу леммы 1) найдутся такие $\nu_n(x)$, для которых 


$$ 


\int_R f(t)d\mu(t)=\lim_{n\to\infty}\int_R\bigg\{ 


\int^\sigma_{-\sigma}


e^{ixt}d\nu_n(x)\bigg\}d\mu(t)=\lim_{n\to\infty} 


\int^\sigma_{-\sigma}\wh\mu(x)d\nu_n(x)=0, 


$$ 


т\,.е. $\cal N^\bot_{\sigma,0}\subset\cal N^\bot_\sigma\subset 


B^\bot_\sigma$, что и завершает доказательство. 


\end{pf} 


 


{\bf 3.} Отметим некоторые свойства этих аннуляторов. 


 


1) Если $\mu(t)\in\cal M^\bot_\sigma$, то $\forall t_0\in R$\; 


$\pm\mu(t\pm t_0)\in\cal M^\bot_\sigma$. Это следует из того, 


что $\pm f(t\pm t_0)\in\cal M_\sigma$, если $f(t)\in\cal M_\sigma$. 


 


2) Если $\mu(t)\in\cal M^\bot_\sigma$ и $\mu(t)\uparrow$, то 


$\wh\mu(x)=0$ при любом $x\in R$, т.\,к. $\forall x\in 


R:|\wh\mu(x)|\le\wh\mu(0)=0$. 


 


3) Если $\mu(t)\in\cal M^\bot_\sigma$ и $\mu(t)\uparrow$, то 


$\mu(t)\equiv\const$. 


 


В самом деле, по формуле обращения $\forall x\in R$ имеем 


$$ 


\mu(x+0)-\mu(x-0)=\lim_{T\to\infty}\frac{1}{2T} 


\int^T_{-T}\wh\mu(t)e^{-ixt}dt=0. 


$$ 


 


Кроме того, в силу монотонности меры $\mu(x)$, 


$\mu(x-0)\le\mu(x)\le\mu(x+0)$, 


т.\,е. $\mu(x)$ непрерывна во всех точках $x\in R$. 


 


Далее, из $\wh\mu(0)=0$ следует, что $\mu(-\infty)=\mu(+\infty)$. 


Отсюда и в силу того, 


что $\forall x\in R:\mu(-\infty)\le\mu(x)\le\mu(+\infty)$, получаем 


$\forall x\in R:\mu(x)=\mu(+\infty)=\const$. 


 


4) Если $\mu(t)=\mu_1(t)+i\mu_2(t)\in\cal M^\bot_\sigma$, 


$\mu_j(t)\uparrow$ ($j=1,2$), то $\mu(t)\equiv\const$. 


 


5) Если $\mu(t)\in\cal M^\bot_\sigma$ и $\sigma\ge\pi$, то она не может 


быть чистой функцией 


скачков в точках $t_k=k$ ($\pm k=0,1,2,\dotsc$), если только функционал 


$\not\equiv0$. 


 


В самом деле, допуская обратное, получим 


$$ 


\wh\mu(x)=\int_R e^{ixt}d\mu(t)=\sum^\infty_{k=-\infty} 


\lambda_k\exp(ikx). 


$$ 


 


Отсюда видно, что $\wh\mu(x)$ является $2\pi$-периодической функцией и 


ряд 


справа сходится абсолютно. Кроме того, $\wh\mu(x)=0$ ($|x|\le\sigma$), 


$\sigma\ge\pi$, поэтому 


$\mu(x)=0$ и во всех точках периода $[-\pi,\pi]$, т.\,е. $\wh\mu(x)=0$ 


$\forall x\in R$. Тогда 


$$ 


\lambda_k=\frac{1}{2\pi}\int^\pi_{-\pi} 


\wh\mu(x)e^{ikx}dx=0\quad (\pm k=0,1,2,\dotsc),


$$ 


но $\lambda_k{=}d_k$ --- скачки функции $\mu(t)$ в точках $t=k$. Значит, 


все скачки 


$\mu(t)$ равны нулю ($d_k{=}\lambda_k{=}0$), а это возможно лишь тогда, 


когда $\mu(t)=\const$, ---


противоречиe условию, что функционал $\not\equiv0$. 


 


6) Если непостоянная функция $\mu(t)\in\cal M^\bot_\sigma$, то не 


существует никакого 


конечного интервала $(a,b)$ ($-\infty<a<b<+\infty$), вне которого 


$\mu(t)=\const$. 


 


В самом деле, если бы существовал такой интервал, то 


$$ 


\wh\mu(x)=\int_R e^{ixt}d\mu(t)=\int^b_a e^{ixt}d\mu(t). 


$$ 


 


Поэтому $\wh\mu(z)$ будет целой функцией конечной степени 


$\le\sigma\le\max(|a|,|b|)$. 


Поскольку $\wh\mu(x)=0$ при $|x|\le\sigma$, то $\wh\mu(z)\equiv0$, 


т.\,е. $\mu(x)=\const$. 


 


7) Если $\mu(t)\in\cal M^\bot_\sigma$ и непостоянна, то какова бы ни 


была $M>0$, существует 


точка роста $t_0\in R$ функции $\mu(t)$, удовлетворяющая условию 


$|t_0|>M$. (Другая формулировка 6).) 


 


8) Если непостоянная $\mu(t)\in\cal M^\bot_\sigma$, то она не может 


иметь конечное число 


точек роста. 


Это утверждение является следствием 6) и 7). 


\medskip 


 


{\bf 4.} На примере покажем, что существует мера $\mu(t)\in\cal 


M^\bot_\sigma$, которая 


сосредоточена на некоторой последовательности $\{t_k\}$, т.\,е. 


является чистой 


функцией скачков с бесконечным числом точек роста. 


 


Пусть $\{\lambda_k\}$ --- некоторая последовательность вещественных 


чисел таких, что 


$\inf\limits_{n\ne m}|\lambda_n-\lambda_m|\ge\delta>0$ и $\psi(z)$ --- 


целая функция, обладающая свойствами 


 


а) $\psi(i\lambda_n)=0$ ($\forall n$), $\psi'(i\lambda_n)\ne0$ 


($\forall n$); 


 


б) вне кружков $\Gamma_n$ с центрами в точках $i\lambda_n$ и радиусами 


$\delta$ 


$$ 


0<m<|\psi(z)|\{e^{\sigma|x|}(1+|x|^\gamma)\}^{-1}<M\ \; (\gamma>1). 


$$ 


 


Положим 


$$ 


\wh\mu(x)=\int_R\exp(ixt)d\mu(t), 


$$ 


где $\mu(t)$ --- чистая функция скачков со скачками 


$d_n=[\psi'(i\lambda_n)]^{-1}$ в точках 


$t_n=\lambda_n$. Тогда будем иметь 


$$ 


\wh\mu(x)=\sum^\infty_{n=-\infty}e^{i\lambda_n x} 


[\psi'(i\lambda_n)]^{-1}. 


$$ 


Остается сослаться на теорему Б.Я.\,Левина ([20], с.\,516), 


в силу которой $\wh\mu(x)=0$ 


при $-\sigma\le x\le\sigma$. 


 


\begin{notenonum} 


Из $\wh\mu_1(x)=\wh\mu_2(x)=0$ при $-\sigma\le x\le\sigma$ 


не следует совпадение $\mu_1(x)$ и $\mu_2(x)$ $\forall x\in R$. 


Действительно, если положить 


$$ 


\wh\mu_1= 


\begin{cases} 


0, & |x|\le\sigma; \\ x^{-2}, & |x|>\sigma, 


\end{cases} 


\qquad \text{и} \qquad 


\wh\mu_2= 


\begin{cases} 


0, & |x|\le\sigma; \\ x^{-4}, & |x|>\sigma \ \; (\sigma>0), 


\end{cases} 


$$ 


то соответствующие им меры $\mu_1(x)$ и $\mu_2(x)$ определяются 


формулами 


$$ 


\mu'_1(x)=\frac{1}{\pi}\int^\infty_\sigma 


\frac{\cos xt}{t^2}dt,\quad 


\mu'_2(x)=\frac{1}{\pi}\int^\infty_\sigma 


\frac{\cos xt}{t^4}dt. 


$$ 


Отсюда видно, что $\mu_1(x)\ne\mu_2(x)$ $\forall x\in R$. 


\end{notenonum} 


 


\section{Соотношения двойственности} 


 


{\bf 1.} Пусть $E_\sigma$ обозначает $\cal M_\sigma$ или $\cal 


N_{\sigma,0}$, $E_{\sigma,1}$ --- единичная сфера в $E_\sigma$, 


$E^\bot_\sigma$ 


--- аннулятор $E_\sigma$, $B_{\sigma,1}$ --- единичная сфера в 


$B_\sigma$. 


 


\begin{thm} 


Пусть $\mu_0(t)$ --- любая мера из $V_R$. Тогда справедливо соотношение 


двойственности 


\begin{equation} 


\sup_{f\in E_{\sigma,1}}\bigg|\int_R f(t)d\mu_0(t)\bigg|= 


\inf_{\mu\in E^\bot_\sigma}\int_R|d\mu_0-d\mu|= 


\sup_{f\in B_{\sigma,1}}\bigg|\int_R f(t)d\mu_0(t)\bigg|, 


\end{equation} 


причем существуют $f_0(t)\in B_{\sigma,1}$ и $\mu^*(t)\in 


E^\bot_\sigma$, которые реализуют $\sup$ и 


$\inf$ в этом соотношении. 


\end{thm} 


 


\begin{pf} 


Первое равенство следует из того, что $E_\sigma$ является 


подпространством $C_0$, а $\mu_0(t)$ определяет некоторый функционал 


$l_0$ из $C^*_0$. 


Докажем последнее равенство. 


 


Пусть $\mu(t)\in E^\bot_\sigma$. В силу теоремы 2\; $\mu(t)\bot 


B_\sigma$. Поэтому имеем 


\begin{gather} 


l_0(f)=\int_R f(t)d\mu_0(t)=\int_R f(t) 


[d\mu_0(t)-d\mu(t)], \notag \\ 


% 


\sup_{f\in B_{\sigma,1}}|l_0(f)|\le 


\inf_{\mu\in E^\bot_\sigma}\int_R 


|d\mu_0(t)-d\mu(t)|. 


\end{gather} 


 


С другой стороны, 


\begin{equation} 


\sup_{f\in B_{\sigma,1}}|l_0(f)|\ge 


\sup_{f\in E_{\sigma,1}}|l_0(f)|=\inf_{f\in E^\bot_{\sigma}} 


\int_R|d\mu_0(t)-d\mu(t)|, 


\end{equation} 


т.\,к. $\cal M_\sigma\subset B_\sigma$. Из (2) и (3) следует (1). 


\end{pf} 


 


Теперь рассмотрим вопрос существования экстремальных элементов. По 


определению точной верхней грани $\exists\{f_n(t)\}\subset 


B_{\sigma,1}$ такая, что 


$$ 


\lambda=\sup_{f\in B_{\sigma,1}}|l_0(f)|= 


\sup_{f\in B_{\sigma,1}}\bigg|\int_R f(t)d\mu_0(t)\bigg|= 


\lim_{n\to\infty}\bigg|\int_R f_n(t)d\mu_0(t)\bigg|. 


$$ 


 


С другой стороны, пространство $B_{\sigma,1}$ замкнутое, выпуклое и 


компактное 


(в смысле равномерной сходимости на каждом конечном отрезке). Поэтому 


можно выделить подпоследовательность $\{f_{n_k}\}\subset\{f_n\}$, 


которая сходится равномерно 


на $R$ к некоторой $f_0(t)\in 


B_{\sigma,1}:\lim\limits_{k\to\infty}f_{n_k}(t)=f_0(t)$. Тогда по 


теореме 


Лебега будем иметь 


$$ 


\lambda=\lim_{k\to\infty}\bigg|\int_R f_{n_k}(t)d\mu_0(t)\bigg|= 


\bigg|\int_R f_0(t)d\mu_0(t)\bigg|. 


$$ 


Остается доказать существование меры $\mu^*(t)\in E^\bot_\sigma$, на 


которой реализуется 


$\inf$ в левой части равенства (1). 


 


Пусть $l_0\in C^*_0$ и $\mu_0(t)\in V_R$ --- мера, соответствующая 


$l_0$ ($\mu_0\leftrightarrow l_0$). 


Тогда для любого функционала $l\in E^\bot_\sigma$ имеем 


\begin{gather} 


|l_0(f)|=|(l_0-l)f|=\bigg|\int_R f(t) 


[d\mu_0(t)-d\mu(t)]\bigg|,\notag \\ 


% 


\|l_0\|\le\inf_{l\in E^\bot_\sigma}\|l_0-l\|= 


\inf_{\mu\in E^\bot_\sigma}\int_R|d\mu_0-d\mu|. 


\end{gather} 


 


С другой стороны, по теореме Хана--Банаха существует функционал $L\in 


C^*_0$, 


для которого выполняются условия 


$$ 


\forall f\in E_\sigma\ \; L(f)=l_0(f)\ \; \text{и} 


\ \; \|L\|=\|l_0\|. 


$$ 


 


Если теперь положим $l^*=L-l_0$, то будем иметь 


$$ 


l^*(f)=0\quad 


(\forall f\in E_\sigma\ \; \text{и}\ \; \forall f\in B_\sigma), 


$$ 


т.\,е. $l^*\in E^\bot_\sigma$. Таким образом, получаем равенства 


$$ 


\|l_0\|=\|L\|=\|l_0-l^*\|=\int_R|d\mu_0-d\mu^*|, 


$$ 


которые показывают, что для функционала $l^*\leftrightarrow\mu^*$ 


неравенство (4) обращается 


в равенство и $l^*$ является экстремальным функционалом (т.\,е. 


$\mu^*\in E^\bot_\sigma$ --- 


экстремальная мера). 


\medskip 


 


{\bf 2. Теорема 4.} {\it Пусть $x(t)\in C_0\setminus B_\sigma$ --- 


любая заданная функция. Тогда 


справедливо соотношение двойственности 


\begin{equation} 


A_\sigma(x;C)=\inf_{y\in B_\sigma}\|x-y\|_C= 


\inf_{y\in E_\sigma}\|x-y\|_C=\sup_{\mu\in E^\bot_{\sigma,1}} 


\bigg|\int_R x(t)d\mu(t)\bigg|, 


\end{equation} 


где $E^\bot_{\sigma,1}$ --- единичная сфера в $E^\bot_\sigma$. 


Существуют $y_0\in B_\sigma$ и $\mu^*\in E^\bot_{\sigma,1}$, на 


которых достигаются $\inf$ и $\sup$ в соотношении $(5)$. 


} 


 


\begin{pf} 


Достаточно показать, что справедливо соотношение 


$$ 


\inf_{y\in B_\sigma}\|x-y\|_C=\sup_{\mu\in E^\bot_{\sigma,1}} 


\bigg|\int_R x(t)d\mu(t)\bigg|. 


$$ 


Для этого заметим, что любая мера $\mu(t)\in E^\bot_\sigma$ 


ортогональна $B_\sigma$, поэтому $\forall y\in B_\sigma$ 


\begin{gather} 


\bigg|\int_R x(t)d\mu(t)\bigg|= 


\bigg|\int_R [x(t)-y(t)]d\mu(t)\bigg|\le 


\|x-y\|_C \int_R|d\mu|,\notag \\ 


% 


\sup_{\mu\in E^\bot_{\sigma,1}} 


\bigg|\int_R x(t)d\mu(t) \bigg|\le\inf_{y\in B_\sigma} 


\|x-y\|_C. 


\end{gather} 


 


С другой стороны, $E_\sigma\subset B_\sigma$, поэтому 


\begin{equation} 


\inf_{y\in B_\sigma}\|x-y\|_C\le\inf_{y\in B_\sigma} 


\|x-y\|_C=\sup_{\mu\in E^\bot_{\sigma,1}} 


\bigg|\int_R x(t)d\mu(t)\bigg|. 


\end{equation} 


Сопоставление (6) и (7) приводит к (5). 


\end{pf} 


 


Существование целой функции $y^*(t)\in B_\sigma$, наименее 


уклоняющейся от 


заданной непрерывной функции, является известным фактом (напр., [3], [4]). 


 


Существование $\mu^*(t)\in E^\bot_{\sigma,1}$ следует из того, что 


$E^\bot_{\sigma,1}$ --- выпуклое и замкнутое 


пространство. Пусть $\{\mu_n\}\subset E^\bot_{\sigma,1}$ и 


$\mu_n\to\mu^*$. Тогда в силу того, что 


$E_\sigma\subset C_0$, имеем $\forall y\in E_\sigma$


$$ 


\int_R y(t)d\mu_n(t)=0,\quad 


\lim_{n\to\infty}\int_R y(t)d\mu_n(t)=0,\quad 


\int_R y(t)d\mu^*(t)=0. 


$$ 


Это означает, что $\mu^*\in E^\bot_\sigma$. Далее, 


$$ 


\lambda=\sup_{\mu\in E^\bot_{\sigma,1}} 


\bigg|\int_R x(t)d\mu(t)\bigg|=\lim_{n\to\infty} 


\bigg|\int_R x(t)d\mu_n(t)\bigg|= 


\bigg|\int_R x(t)d\mu^*(t)\bigg|, 


$$ 


т.\,к. $x(t)\in C_0.\quad\square$ 


 


{\bf 3.} В конце заметим, что выбором меры $\mu_0(t)$ можно решать


задачи о нахождении таких величин, как 


$$ 


\max_f\bigg|f\bigg(x_0+\frac{\pi}{2\sigma} \bigg)-f 


\bigg(x_0-\frac{\pi}{2\sigma} \bigg)\bigg|,\qquad 


\max_f\bigg|\sum^n_{k=1}\lambda_k f(x_k)\bigg| 


$$ 


при заданных $\lambda_k$ и $x_k$; 


$$ 


\max_f|f'(x_0)|,\quad \max_f|Af(x_1)+B f'(x_2)|,\quad 


\max_f|\sin\alpha f'(x)-\sigma\cos\alpha f(x)| 


\ \; \text{и т.\,п.,} 


$$ 


которые были предметом специальных исследований. 


 


Например, если $\mu_0(t)$ будет функцией скачков со скачками 


$d_k=\frac{4\sigma}{\pi^2}\frac{(-1)^{k-1}}{(2k+1)^2}$ 


в точках $t_k=x_0+\frac{(2k+1)\pi}{2\sigma}$ ($\pm k=0,1,2,\dotsc$), 


$x_0\in R$, то 


$$ 


l_0(f)=\frac{4\sigma}{\pi^2}\sum^\infty_{k=-\infty} 


\frac{(-1)^{k-1}}{(2k+1)^2}f\bigg(x_0+\frac{2k+1}{2\sigma}\pi \bigg) 


\equiv f'(x_0). 


$$ 


 


Тем самым получена задача С.Н.\,Бернштейна о нахождении 


$\max\limits_{f\in B_{\sigma,1}}|f'(x_0)|$ 


([2]; [4], с.\,74) и ее решение 


$$ 


\max_{f\in B_{\sigma,1}}|f'(x_0)|=\sigma. 


$$ 


 


\begin{thebibliography}{99} 


 


\bibitem{1} 


Ахиезер Н.И. {\em Лекции по теории аппроксимации}. -- 2-е изд. -- М.: 


Наука, 1965. -- 407 с. 


\bibitem{2} 


Bernstein S.N. Sur une propri\'et\'e des fonctions enti\`eres // 


Compt. Rend. Acad. Sci. -- 1923. -- \No\,176. -- P.\,1603--1605. 


\bibitem{3} 


Тиман А.Ф. {\em Теория приближения функций действительного 


переменного}. 


-- М.: Физматгиз, 1960. -- 624 с. 


\bibitem{4} 


Ибрагимов И.И. {\em Экстремальные свойства целых функций конечной 


степени}. -- Баку: Изд.-во АН Азерб\,ССР, 1962. -- 373 с. 


\bibitem{5} 


Ибрагимов И.И. {\em Теория приближения целыми функциями}. -- Баку: 


Изд-во ``Элм'', 1979. -- 468 с. 


\bibitem{6} 


Насибов Ф.Г. {\em Об экстремальных задачах в классе $B_\sigma$} // 


ДАН СССР. -- 1987. -- T.\,294. -- \No\,2. -- C.\,268--271. 


\bibitem{7} 


Насибов Ф.Г. {\em Аннулятор класса $W_{\sigma,2}$ и некоторые его 


применения}. -- Азер. гос. нефт. акад. -- Баку, 1982. -- 11 c. --


Деп. в ВИНИТИ, \No\,1671-82. 


\bibitem{8} 


Насибов Ф.Г. {\em O приближении в $L_2$ целыми функциями} // 


ДАН Азерб\,ССР. -- 1986. -- T.\,42. -- \No\,4. -- C.\,3--6. 


\bibitem{9} 


Насибов Ф.Г. {\em О приближении функций в равномерной метрике и в 


среднем} // Исследов. по некоторым вопр. конструктивной теории 


функций 


и дифференц. уравнений. -- Баку: Аз.\,ИНЕФТЕХИМ, 1983. -- 


C.\,23--31. 


\bibitem{10} 


Насибов Ф.Г. {\em Описание аннулятора одного класса целых функций 


конечной полустепени и некоторые приложения} // 


Изв. вузов. Математика. -- 1986. -- \No\,2. -- C.\,29--33. 


\bibitem{11} 


Хавинсон С.Я. {\em Об одной экстремальной задаче теории аналитических 


функций} // УМН. -- 1949. -- Т.\,IV. -- Bып.\,4. -- C.\,158--159. 


\bibitem{12} 


Хавинсон С.Я. {\em Экстремальные задачи для некоторых классов 


аналитических функций в конечносвязных областях} // 


Mатем. сб. -- 1955. -- \No\,36. -- C.\,445--478. 


\bibitem{13} 


Хавинсон С.Я. {\em Теория экстремальных задач для ограниченных 


аналитических функций, удовлетворяющих дополнительным условиям внутри 


области} // УМН. -- 1963. -- Т.\,XVIII. -- Bып.\,2. -- C.\,23--98. 


\bibitem{14} 


Тумаркин Г.Ц., Хавинсон С.Я. {\em Исследование свойств экстремальных 


функций с помощью соотношений двойственности в экстремальных 


задачах для классов аналитических функций в многосвязных областях} // 


Матем. сб. -- 1959. -- Т.\,46. -- Bып.\,2. -- C.\,195--228. 


\bibitem{15} 


Ахиезер Н.И., Крейн М.Г. {\em О некоторых вопросах теории моментов}. -- 


Харьков: ОНТИ, 1938. -- 253 с. 


\bibitem{16} 


Никольский С.М. {\em Приближение функций тригонометрическими полиномами 


в среднем} // Изв. АН СССР. Сер. матем. -- 1946. -- \No\,10. -- 


C.\,207--256. 


\bibitem{17} 


Горин Е.А. {\em Неравенства Бернштейна с точки зрения теории 


операторов} 


// Вестн. Харьковск. ун-та. -- 1980. -- \No\,205. -- C.\,77--105. 


\bibitem{18} 


Лукач Е. {\em Характеристические функции}. -- М.: Наука, 1979. -- 423 


с. 


\bibitem{19} 


Boas R. {\em Entire functions}. -- New~York: Acad. Press, 1954. -- 276 


p. 


\bibitem{20} 


Левин Б.Я. {\em Распределение корней целых функций}. -- М.: 


Гостехиздат, 


1956. -- 632 с. 


 


\end{thebibliography} 


 


\vskip 2cm 


 


\postup{\it Азербайджанская государственная}{$\qquad$\it Поступили} 


\postup{\it нефтяная aкадемия}{\it первый вариант $10.06.1997$} 


\postup{}{\it окончательный вариант $16.08.2000$} 


 


\label{end} 


\end{document} 


